\documentclass[a4paper,12pt]{article}
\usepackage[utf8]{inputenc}
\usepackage[T2A]{fontenc}
\usepackage[russian]{babel}
\usepackage{hyperref}
\usepackage{graphicx}
\usepackage{geometry}
\usepackage{array}

\geometry{top=2cm, bottom=2cm, left=2.5cm, right=2.5cm}
\hypersetup{
    colorlinks=true,
    linkcolor=blue,
    urlcolor=blue,
}

\begin{document}

\begin{titlepage}
    \begin{center}
        \textbf{\Large Регламент практик по ООП 2025-2026 \\ (весенний семестр)}
        \vspace{5cm}
    \end{center}
\end{titlepage}

\newpage
\tableofcontents
\newpage

\section{Введение}

Весенний семестр курса ООП посвящен изучению универсальных технологий программирования применительно к языку Java.

Цели студентов на практике по технологиям программирования:
\begin{enumerate}
    \item Научиться применять программирование в объектно-ориентированном стиле для декомпозиции сложной системы. Используя обратную связь от преподавателя, улучшить умение проектирования сложных систем с помощью приемов ООП.
    \item Получить практический навык в следующих областях: построение программ с графическим интерфейсом, многопоточное программирование, загрузка классов, сериализация классов, сетевое взаимодействие, управление памятью.
\end{enumerate}

\section{Задачи}

Крайне редко в реальной жизни проект начинается с нуля, большую часть времени программист работает с уже существующим кодом, развивает существующую систему. Поэтому в рамках этого семестра моделируем подобную ситуацию.

За основу берем существующий проект, размещенный здесь: \\
\url{https://github.com/valeevr/base-robots}

В коде проекта есть как важные для изучения моменты (как организовать многооконное приложение с графическим интерфейсом), так и ошибки с точки зрения проектирования, мешающие дальнейшему развитию (вполне реальная ситуация). Поэтому сначала следует разобраться с имеющимся кодом и попробовать устранить проблемы (см. стартовые задачи), а затем начать развивать проект.

Код решения задач должен быть покрыт модульными тестами. Количество и качество тестов влияют на оценку задачи.

Для каждой задачи, которую предполагается решать, совместно с преподавателем устанавливается крайний срок сдачи (каждая задача должна решаться в срок 3 недели с момента ее взятия в работу). После истечения этого срока задача обесценивается (теряет в стоимости 50\%).

\section{Стартовые задачи}

\subsection{Задача 1: Знакомство с устройством MDI приложения на Java (8 баллов)}

В предоставленном комплекте исходников уже реализована некоторая функциональность MDI (multi-document interface) приложения (главное окно, меню, окно протокола, окно моделирования движения робота). Требуется разобраться со следующими моментами:
\begin{itemize}
    \item как запускается приложение?
    \item как создаются окна?
    \item как инициализируется меню?
    \item как подключаются обработчики событий?
\end{itemize}

\textbf{Задания:}
\begin{itemize}
    \item требуется добавить пункт меню, позволяющий закрыть приложение;
    \item требуется собрать обработку события выхода из приложения в один метод и сделать так, чтобы в этом методе выдавался запрос на подтверждение выхода (см. класс JOptionPane);
    \item сделать так, чтобы диалог на основе JOptionPane выдавал текст на кнопках на русском языке.
\end{itemize}

\textbf{Дополнительное задание:}
\begin{itemize}
    \item попробовать разобраться, как все стандартные предопределенные компоненты заставить работать на русском языке (+10 баллов).
\end{itemize}

\subsection{Задача 2: Знакомство с управлением разнородными объектами (16 баллов)}

Обычно считается хорошим тоном, когда приложение запоминает свое состояние на момент выхода, чтобы можно было продолжить работу с ним без дополнительной настройки при последующем запуске.

\textbf{Задания:}
\begin{itemize}
    \item в методе выхода из приложения нужно сохранить (в файл) положения окошек (положение на экране, развернутое/свернутое состояние);
    \item при запуске приложения следует восстановить геометрию окошек приложения (положение на экране, развернутое/свернутое состояние);
\end{itemize}
Файл с конфигурацией следует создавать в домашнем каталоге пользователя.

\subsection{Задача 3: Знакомство с MVC (16 баллов)}

Можно заметить, что при определенных условиях робот ведет себя странно — уползает за пределы окна, и что с ним происходит дальше — непонятно.

Надо исправить ошибку, и для этого нужно добавить в программу новое окно, в котором отображались бы координаты робота при его движении.

\textbf{Задания:}
\begin{itemize}
    \item реализовать новое окно (например, на основе диалоговых компонент), в котором будут отображаться текущие координаты робота;
    \item для этого придется разделить логику отображения робота (визуализацию) и модель робота (код получения координат и моделирования движения);
    \item модель нужно будет описать как обозреваемый объект (Observable), чтобы затем окно могло подписаться на уведомления от модели и обновлять свое состояние;
\end{itemize}

\section{Последующие задачи}

После успешного решения стартовых задач студенты сами выбирают новые возможности, которые собираются добавить в приложение. Всю работу над проектом студенты делят на задачи. Преподаватели оценивают каждую задачу независимо.

Задача — это объем работ с четко сформулированной одной целостной целью, которую можно объяснить за 2 минуты. Формулировка задачи должна быть зафиксирована письменно (можно использовать возможности по описанию и ведению задач в сервисах GitHub, GitLab или подобных) и не требовать дополнительных пояснений.

Задача должна затрагивать материал, изучаемый на лекционных занятиях. А именно, должна касаться тем из следующего списка:
\begin{itemize}
    \item нетривиальные возможности при построении графического интерфейса;
    \item локализация программ;
    \item отчуждение данных, сериализация, версионирование;
    \item многопоточное программирование;
    \item использование синглтонов, кэширования, конвейеров;
    \item загрузка классов;
    \item использование ссылок, отличных от жестких;
    \item управление владением;
    \item асинхронная работа с сетью;
    \item использование рефлексии.
\end{itemize}
Применение каждой конкретной техники должно быть оправданным!

Рекомендованный ориентировочный объем работы на одну задачу — 200 строк рабочего кода, 2–5 классов, модульные тесты, ~ 4 часа работы.

Студенты предлагают свои формулировки задачи. Преподаватель корректирует формулировку, чтобы сделать задачу более полезной с точки зрения обучения: предлагает дополнительные требования, чтобы пришлось применить приёмы проектирования, или советует, как сократить объем работы для студентов. Здесь же можно оценить сложность задачи в баллах.

Сдать задачу можно только после того, как преподаватель явно согласился с её формулировкой. Валидировать формулировки задач лучше заранее (до предполагаемого начала работы над задачей) через чат  с преподавателем.

\begin{center}
\begin{tabular}{|p{10cm}|}
    \hline
    Сдать задачу можно только после того, как преподаватель явно согласился с её формулировкой. Валидировать формулировки задач лучше заранее (до предполагаемого начала работы над задачей) через чат  с преподавателем. \\
    \hline
\end{tabular}
\end{center}


\section{Возможные идеи задач}

\subsection{Пример 1: Нетривиальная структура (10+10 баллов)}

Придумать и реализовать структуру данных для хранения записей, отображаемых в окне протоколирования. В исходной версии кода в качестве такой структуры использован ArrayList без какой-либо самоочистки. В реальности это приведет к ненужному накоплению информации в логе (и такой своеобразной утечки памяти, необходимой для хранения старых записей лога). Поэтому обычно требуется, чтобы размер хранимой (и отображаемой в окне) части протокола был ограничен. Но если просто удалять из ArrayList старые записи, это приведет к неэффективной реализации (добавление записи будет работать за O(N), поскольку придется удалять устаревшие записи из начала структуры) и к проблемам с конкурентным доступом (возникает состояние гонки между добавлением данных и итерированием существующих).

Соответственно, структура для хранения протокола должна обладать следующими свойствами:
\begin{itemize}
    \item имела ограниченный размер (старые записи вытесняются);
    \item должна быть потокобезопасной (запись и чтение порождают состояние гонки);
    \item должна быть возможность доступа к части данных (сегмент смежных записей) по индексам начала и конца (такая операция, по идее, нужна для эффективного отображения данных в окне, чтобы не читать полный лог);
    \item добавление данных (с потенциальным удалением старых) и чтение уже хранящихся должны быть быстрыми (т.е. никаких O(n));
    \item должна возвращать потокобезопасный итератор (то есть добавление данных во время итерирования не должно разрушать итератор).
\end{itemize}

Первые 3 пункта могут быть реализованы на любой структуре (стэк, очередь, массив), но потребуется сделать структуру потокобезопасной (конкурентные чтение и запись).

Последние 2 пункта — усложненная версия задачи, позволяющая познакомиться, например, с темпоральными структурами данных.

\subsection{Пример 2: Локализация (10+10 баллов)}

Сделать так, чтобы приложение было локализовано (то есть обеспечить возможность выбора через меню или диалог настроек языка, на котором работает пользовательский интерфейс).

Для упрощения задачи можно сделать следующие допущения:
\begin{enumerate}
    \item локализацию проводим только для собственных компонент программы (сообщения в готовых диалогах Swing переводить не нужно);
    \item считаем, что должны поддерживаться два языка: русский (выводится в обычном русском алфавите) и, на выбор, транслит (русский язык, выводимый латинницей) или английский язык (если нет сложностей с переводом).
\end{enumerate}

Для решения задачи следует разобраться с работой классов ResourceBundle и с тем, как формировать сообщения через форматирование.

\textbf{Усложнение задачи:}
Ускорить форматирование сообщений за счет использования кэша. Здесь важно разобраться с тем, что Java предоставляет два механизма для шаблонизации: на основе класса MessageFormat, и на основе класса Formatter. В случае с классом MessageFormat можно заметить, что вызов MessageFormat.format сводится к запуску парсера (создает объект MessageFormat), который потом используется для форматирования. Так вот в кэше можно сохранять эти самые созданные объекты MessageFormat, и тогда форматирование будет работать существенно быстрее (парсер работает дольше, чем код форматирования). В случае с классом Formatter такая техника, увы, не работает. После реализации следует провести сравнительный анализ скорости форматирования с помощью класса Formatter, с помощью класса MessageFormat без кэширования и с помощью класса Formatter с кэшированием.

\subsection{Пример 3: Разные роботы (20 баллов)}

Сделать так, чтобы можно было загружать "роботов" извне. А именно, добавить возможность выбрать файл .jar, в котором были бы собраны классы (файлы с байт-кодом), содержащие логику работы робота (например, робот может быть 4-колесным) и правила его визуализации (картинка робота, правила его отображения).

Задача направлена на то, чтобы получить возможность поработать с загрузчиками классов. При загрузке файла следует убедиться, что это именно jar-архив, а затем создать URLClassLoader, чтобы загрузить необходимые классы.

\section{Важные дополнительные изменения}

\begin{enumerate}
    \item задачи из примеров этого регламента должны чередоваться с заданиями с персональными формулировками
    \item есть дополнительная возможность заработать до 20 баллов в семестре за активности на парах
\end{enumerate}


\section{Теория}

Если для выполнения задачи студентам требуются знания, ещё не пройденные на лекциях, студенты самостоятельно осваивают эту тему по лекционным слайдам или по аудиоверсии курса. Они доступны по ссылке: \\
\url{http://courses.imkn.urfu.ru/oop/}

\section{Работа в командах}

Работу над проектом можно строить как индивидуально, так и в команде (в последнем случае можно успеть сделать больше и получить более интересный результат). В команде может быть до 3 человек. Командная работа допускается только после решения каждым из участников какой-то из стартовых задач, причем разными участниками команды должны быть решены разные стартовые задачи (репозитарии затем можно и нужно объединить).

Результат работы одного студента в команде должен проходить ревью кем-то другим из команды. Каждый в паре должен нести ответственность за весь код задачи, уметь объяснить и защитить каждое принятое решение.

Есть несколько шаблонов работы в паре (в команде). Выбирайте!
\begin{enumerate}
    \item Вы весь код пишете в паре, собравшись в одном месте физически. Пишете одновременно код и модульные тесты на него. Так оба участника будут знать всё про задачу. Эта практика называется парным программированием или экстремальным программированием.
    \item Вы декомпозируете задачу на подзадачи и каждый занимается своей частью. Каждый пишет и код, и тесты для своей части решения. Затем партнеры по команде выполняют взаимное ревью кода, добавляют тесты. В конце код сливается в общее решение.
\end{enumerate}

Следует отметить, что нельзя делиться по принципу: один человек только пишет код, а второй — только тестирует. В реальной работе разработчик всегда пишет модульные тесты на свой код самостоятельно, в идеале параллельно с написанием кода (TDD, Test Driven Development).

В любой момент студенты могут поменять пару. В любой момент преподаватель может переформировать пары другим способом. Новая пара студентов выбирает одну из двух кодовых баз, создает её копию (fork на языке github) и продолжает работать над ней. Может быть так, что две пары после обмена участниками продолжат работу над одной и той же кодовой базой. Это нормально.

Распределение по парам согласуется с преподавателем. Детали обсуждаются в чате в Телеграм.

За счет работы в команде можно набрать на 30\% больше баллов, если при работе в команде действительно наблюдается активное сотрудничество членов команды (активное взаимное ревью, активное вовлечение в понимание особенностей задач, решаемых другими участниками команды и т.п.).

\section{Git}

Весь исходный код должен быть в системе контроля версий git. Рекомендованная система — \url{https://gitlab.com}, но разрешено использовать и аналоги, например \url{https://github.com}

Как освоиться с гитом? От простых способов к сложным:
\begin{enumerate}
    \item Пройти краткий гайд.
    \item Пройти интерактивные учебные курсы от github и schoolacademy.
    \item Прочитать официальную книгу по git: \url{http://git-scm.com/book/ru/v2} Первые три главы обязательны для уверенного использования git.
\end{enumerate}

Снизить количество ошибок при использовании гита можно, если для операций коммита и просмотра изменений использовать один из графических клиентов. Подойдут, как отдельные приложения (Git Extensions, TortoiseGit (только для Windows), GitHub Desktop или SourceTree, веб-инструменты GitHub и GitLab для выполнения слияния веток), так и встроенные в среды разработки (в IntelliJ IDEA есть встроенные инструменты для работы с Git).

Для того, чтобы команде было комфортно работать с общим исходным кодом важно соблюдать гигиенические требования при пользовании системы контроля версий:
\begin{enumerate}
    \item Git-репозиторий должен быть самодостаточным. То есть любой разработчик, работающий с вашим проектом должен иметь возможность скачать исходники в новый чистый каталог, скомпилировать и запустить проект без необходимости решать проблемы с отсутствующими файлами. В частности, в проекте не должны использоваться файлы, присутствующие только на вашей рабочей машине.
    \item Git-репозиторий не должен содержать ничего лишнего. Файлы, которые появляются в процессе компиляции, файлы с пользовательскими настройками вашей IDE, бинарные зависимости (если они получаются менеджером пакетов, например, maven) — все это должно быть включено в .gitignore-файл, и не должно попадать в репозиторий.
    \item Текстовые сообщения к коммитам должны быть содержательными (то есть по их тексту должно быть ясно, что именно сделано в данном коммите). Также рекомендуется в текст комментария включать какую-то метку задачи (например, в формате Task-NNN, где NNN — это номер задачи), чтобы иметь возможность находить все коммиты, относящиеся к конкретной задаче.
    \item Перед коммитом просматривайте все те изменения, которые вы совершили и не включайте в коммит лишние изменения. Помните, что каждый коммит должен быть самодостаточным (т.е. не следует в один коммит включать несколько изменений кода для разных подзадач). Если так получилось, что накопились правки по нескольким задачам, то можно или сделать серию коммитов, выбирая нужные файлы с изменениями (при этом не забываем, что после каждого коммита код обычно должен компилироваться), или можно создать другую ветку, в которую...
    \item Перед каждым коммитом проверяйте работоспособность всего проекта. Обычно предполагается, что любое состояние проекта в Git по крайней мере компилируется без ошибок.
    \item В git-репозитории не должно быть никаких секретных ключей или паролей. Ибо все, что попадает в облако, следует считать общедоступным (то есть и пароли, и ключи) — конечно же такого никто не хочет.
    \item В git должна быть видна работа над задачами всех участников команды.
    \item При работе над задачами следует использовать ветки. То есть все наработки по задаче выполняются в ветке. Затем выполняется ревью с серией замечаний и исправлений (его можно выполнять в режиме Merge Request (запрос на вливание в master), или в режиме Pull Request (запрос на вливание в чужой репозитарий). И только потом правки (проверенные и стабильные!) можно вливать в основную ветку (master).
\end{enumerate}

Формулировки задач и комментарии к ним также стоит сохранять в git. Это удобно делать или через инструмент по управлению задачами (называется issue tracker), доступный в актуальных версиях всех облачных систем, основанных на Git (тогда в п.3 в качестве идентификатора задач удобно использовать номер задачи в этом самом инструменте по управлению задачами), или в файле Readme.md в корне репозитория (содержимое этого файла показывается на главной странице вашего репозитория). Во всех случаях редактировать задачи или Readme.md можно прямо через браузер.

Код, демонстрируемый преподавателю должен быть закоммичен, т.е. полностью загружен в облачный Git. Нельзя сдавать код, который ещё не находится в репозитории.

За правильное использование системы контроля версий (см. гигиенические требования) можно получить до 3 дополнительных баллов к каждой задаче. За нарушения гигиенических требований можно потерять баллы.

\section{Баллы}

Каждая задача в случае успешной сдачи оценивается в согласованный на этапе формулировки балл.

При работе в команде сумма баллов, полученных за задачи, делится на количество участников команды.

На каждом занятии можно попытаться сдать только одну-две задачи (из них должна быть лишь одна новая). Сдавать задачу можно только после того, как преподаватель принял вашу формулировку задачи или скорректировал её. Ожидайте, что с первого раза сдать задачу не получится — будут претензии к качеству и доработки, которые вы сможете сдать только на второй раз.

Одну задачу нельзя сдавать более трех раз. После третьей попытки её сдать преподаватель ставит либо 0, либо менее половины стоимости задачи по своему усмотрению.

Рекомендуется готовить формулировку следующей задачи заранее и согласовывать её с преподавателем либо через чат группы до пары, либо сразу после первой попытки сдачи предыдущей задачи.

\subsection*{Границы баллов:}
\begin{enumerate}
    \item Уровень «отлично» — 80 баллов и выше.
    \item Для допуска к экзамену достаточно 40 баллов.
    \item Предварительная граница для упрощённой сдачи экзамена: 80 баллов.
\end{enumerate}

\end{document}